\documentclass[11pt,a4paper]{article}
\usepackage[T1]{fontenc}
\usepackage{lmodern,amsmath,amssymb,booktabs,tabularx,microtype}
\usepackage[margin=24mm]{geometry}
\usepackage[bookmarks=false,colorlinks=true,linkcolor=blue,urlcolor=blue]{hyperref}
\setlength{\parindent}{0pt}\setlength{\parskip}{5pt}
\setlength{\emergencystretch}{2em}
\title{BundleFuse: From Prior Work to the Research Question}
\author{}\date{5 November 2026}
\begin{document}
\maketitle\vspace{-18pt}
This document develops the motivation and positioning before the experimental case.
Its purpose is to explain what is inherited, what question is being asked, and how
the implemented construction addresses that question. It is a focused positioning
argument, not an exhaustive literature review or a claim of established priority.
The accompanying technical note defines the exact scores and their assumptions.

\section{Begin with the kind of relationship the task requires}
A partial bundle specifies a context for choosing another item. A camera bag can
complete a camera kit even though its appearance and function differ from the camera.
A second camera may be more similar but less useful as a completion. Thus a useful
ranking must represent the relationship required by the context, rather than assume
that resemblance always implies suitability.

This distinction has an established basis. McAuley, Pandey and Leskovec model
substitutable and complementary product links using product-associated information,
including review text and category structure~\cite{mcauley}. Their work motivates
relation-sensitive recommendation. It does not establish that every pair from
different categories is complementary, or that a category rule solves compatibility.

For bundle completion, the relationship can also vary by domain. Repetition of a
product role can be undesirable in an outfit, while a shared artist can be useful
in a playlist. This motivates retaining several precisely defined relationship
scores. It does not justify assuming that all of them are useful everywhere.

This distinction is already central to fashion compatibility research. Vasileva
et al. learn type-aware embeddings that jointly represent similarity and
compatibility~\cite{vasileva}. Their method learns a representation from images;
BundleFuse instead computes category and descriptive scores. The difference is in
the construction, not the recognition that item type matters.

\section{What explicit recommendation models already establish}
Item-relation methods provide a natural starting point: score a candidate by its
relationships to the observed items. Let $\mathcal I$ be a catalogue, $S$ its observed
subset, $c\notin S$ a candidate, and $M_{cs}$ an estimated relation between $c$ and
$s$. The additive score is
\[
 r(c\mid S)=\sum_{s\in S}M_{cs}.
\]
This notation describes a class of estimators; it is not presented as a new formula.

Sarwar et al. establish the item-neighborhood approach: estimate item-item
relations from user ratings and use them to form predictions~\cite{sarwar}.
BundleFuse applies a related principle to bundle co-memberships, which record
presence rather than explicit ratings. The analogy does not make the two inputs
or their meanings identical.

EASE demonstrates that a shallow linear recommendation model with a closed-form
solution can achieve strong ranking performance~\cite{ease}. Side-information
extensions of SLIM already combine purchase histories and item attributes in a
regularized linear model~\cite{sslim}. More recently, the FEASE preprint extends EASE
with user and item side features for cold-start handling~\cite{fease}.

These works constrain our claim. Explicit scoring, side-information fusion and the
absence of a deep downstream encoder are established possibilities. Furthermore,
a model originally formulated for user histories can often be applied to an observed
item set. Merely renaming a history as a partial bundle does not create a new method.

The remaining design issue for this work is more specific: \emph{which relations
should an explicit estimator preserve when completing bundles, and how should their
scores be combined?} Historical overlap, category relations and descriptive resemblance
are related measurements, but they are not interchangeable definitions of compatibility.

\section{What modern bundle construction already addresses}
CLHE directly addresses sparse bundle histories by combining multimodal features,
item-level user feedback and bundle composition. It uses hierarchical self-attention
and item- and bundle-level contrastive objectives~\cite{clhe}. Its motivation already
includes the observation that features can be available for items with little bundle
history. BundleFuse must acknowledge that motivation as prior work.

RaMen develops characteristic and collaborative strategies together with hypergraph
structure to learn bundle representations~\cite{ramen}. It explicitly models multiple
kinds of relationships. A claim that earlier bundle models consider only similarity,
ignore direct item associations, or neglect multi-item structure would misrepresent
this comparator.

These papers motivate a question about the \emph{construction of the predictor}.
They do not imply that another learned encoder is the only way to use the available
information, nor that their representation learning is unnecessary. That issue is
comparative and must be investigated.

\section{State the question without inventing an unaddressed problem}
The research question is:
\begin{quote}
Given historical bundles and available item attributes, how far can explicit,
relation-specific compatibility estimates support bundle completion without training
a new neural bundle encoder?
\end{quote}
This asks about the usefulness and limits of a particular alternative construction.
It does not assert that no earlier model combines these inputs. Its value depends
on defining a nontrivial construction, explaining its choices, and eventually
showing what those choices achieve relative to appropriate alternatives.

The distinction from neural bundle encoders is that the relevant relationships enter
as directly computed candidate scores. The distinction from a generic linear
recommender must be established through the specific relation constructions and
combination rule, not merely through the claim that the method is simple.

The combination of history and descriptions also has a long history. Melville
et al. use content predictions to enrich sparse user ratings before collaborative
filtering~\cite{melville}. BundleFuse forms several candidate scores from bundle
history and item attributes and then combines them; it does not first fill a user
rating matrix. Both inherit the idea that content can help where collaborative
observations are sparse, so this idea is not itself our contribution.

An even closer task-level precedent is Volkovs et al.'s playlist-continuation
system~\cite{volkovs}. Its retrieval stage blends item- and playlist-neighbor
scores with other collaborative models, adjusts scores for popularity, and is
followed by a learned reranker. BundleFuse's particular relations, scaling and
fallback differ, but combining multiple retrieval signals and correcting
popularity in a partial-set task are already established. A fair comparison must
account for task protocol, inputs, and the learned reranking stage.

\section{Develop the construction from that question}
\subsection{Preserve directly observed relations}
Let $B$ be the binary training-bundle by item matrix and let $x_S$ be the indicator
vector of $S$. Then $(B^\top Bx_S)_c$ counts recorded co-memberships with the query
items. BundleFuse uses weighted overlap and alternative summaries of history to
reduce particular effects of frequency, bundle size and sparse direct observations.
This step inherits the explicit item-relation perspective. Its particular weights
and summary operators are modeling choices, not consequences of the cited papers.

\subsection{Share evidence while retaining the type of relationship}
Exact co-occurrence has no direct evidence for an unobserved pair. BundleFuse also
constructs category relations and scores from descriptive features. Category
relations aggregate observations across items, while feature relationships extend
coverage beyond recorded pairs. The category component can distinguish support
between categories from penalties for repeating a rarely repeated category.

These choices connect the relation-sensitive view to an explicit score construction.
They do not identify economic substitutes or complements from category labels alone.
The distinction is operational: each formula specifies which observed relationship
it measures and what response it assigns to the candidate.

Category or cluster centering of descriptors and graph-adjusted features provide
further candidate estimators. Their purpose is to separate a group's mean descriptor
from within-group variation and to combine descriptions with historical relationships.
These operations are not proofs of disentanglement. Their individual utility remains
an empirical question. The complete method should disclose each operation, even when
several share one conceptual role.

\subsection{Combine scores with different scales and coverage}
Let $r_h(c\mid S)$ denote raw channel $h$, let $z_h(c\mid S)$ be its catalogue-standardized
value for the query, and let $w_h\geq0$ be its validation-selected weight. The additive
core is
\[
 Y_0(c\mid S)=\sum_h w_h z_h(c\mid S).
\]
This exposes the construction to inspection: the meaning of a channel, its scale
transformation and its contribution to a candidate score are explicit. It does not
make channels independent or their weights causally interpretable. The use of
validation to choose weights is fitting, even without neural gradient training.

A closer interpretation of the history component comes from set-neighbor
recommendation. For binary training bundles $b$, define $m_b=|b\cap S|$,
$a_b=\sum_{s\in b\cap S}(f_s+5)^{-\alpha}$ and
$v_b=(|b|+1)^{-\gamma}$. Before scale alignment, the pair score plus its
incidence correction can be written, for nonnegative $\beta$, as
\[
 \sum_{b:c\in b} v_b a_b\bigl[1+\beta\max(0,m_b-1)\bigr].
\]
It is therefore a sum over candidate-containing historical bundles with a
nonlinear query-overlap weight. Playlist- and session-neighbor methods already
score candidates through similar historical sets~\cite{bonnin,ludewig}. Our
particular weighting does not impose a nearest-neighbor cutoff, and its formula
is not claimed to be identical to theirs. Weighted-incidence dataset profiles
need their additional bundle weights retained explicitly.

Further prior art limits a general novelty claim: normalized hypergraph
incidence operators predate this method~\cite{zhou}; playlist hybrids have fused
co-occurrence with audio and text~\cite{ferraro}. Earlier basket-completion models based on determinantal point processes also
model set quality and diversity~\cite{gartrell,warlop}. A detailed map of these and
other overlaps, including unresolved exact-formula questions, is provided in
\texttt{method\_prior\_art\_review.md}.

The signal menu therefore belongs inside the explanation of these relation estimates.
It is not a list of independent novelty claims. A convincing presentation must
explain why each family is considered, state its overlap with other families, and
identify which components survive selection in the reported model.

\subsection{Introduce optional corrections after the core}
Let $Y$ be the combined score including any selected multi-item corrections, and
$\mu_0$ a scaled fallback assembled from a restricted set of channels. Let $f_c$ be
the number of candidate appearances in training and $\kappa>0$ the fallback strength.
The backoff score is
\[
 \theta(c\mid S)=\frac{f_c}{f_c+\kappa}Y(c\mid S)
              +\frac{\kappa}{f_c+\kappa}\mu_0(c\mid S).
\]
It provides an explicit rule for changing the balance with candidate history.
Frequency is a proxy for available historical evidence, not a calibrated uncertainty.
The two estimates overlap; the fallback is not a history-independent expert.

Related empirical-Bayes product-search work already combines behavioral information
with non-behavioral priors~\cite{han}. Consequently, neither interpolation nor
shrinkage toward descriptive evidence is new by itself. BundleFuse's specific
fallback, its overlap with the combined score, and the candidate-frequency rule
must be defined and defended in the bundle setting. The technical note gives a
pointwise regularization interpretation without assuming a validated noise model.
Hu et al. already use an item-statistics-dependent weight to balance behavioral
and content representations~\cite{hu}. CoHeat combines two user--bundle scores
with a bundle-popularity-dependent weight and proves the corresponding
monotonic influence properties~\cite{coheat}. The task, scores and weighting
function differ from BundleFuse, but the general adaptive-mixture principle and
monotonic sensitivity argument are established prior work.

The incidence correction also admits a precise relation to earlier itemset
methods. In the binary, unit-bundle-weight case, put
$d_s=(f_s+5)^{-\alpha}$ and $v_b=(|b|+1)^{-\gamma}$. Expanding its overlap
factors gives
\[
 \Delta(c\mid S)=\beta\!\sum_{\substack{s,t\in S\\s\ne t}}\!d_s
       \sum_b v_b B_{bc}B_{bs}B_{bt}.
\]
Thus the correction aggregates weighted historical co-occurrences of each
seed pair with the candidate. Higher-order neighborhood models and
HOSLIM already score candidates using itemset-to-item associations~\cite{hoslim}.
HOSLIM learns sparse coefficients; this construction uses weighted counts and
specific shared weights. The formula does not imply the two predictors are equal.
The broader idea of amplifying strongly matching neighbors also predates this
construction~\cite{breese,kelen}.

The incidence surplus and Rao-based term address an additional question: can support
involving several observed items change the assessment beyond separate matches?
They are explicit optional corrections. Their presence does not establish that
higher-order interactions are the principal explanation of the method's accuracy.

Higher-order item interactions in recommendation also have prior work: DeepICF
models interactions among previously consumed items with neural
networks~\cite{deepicf}. This is a different mechanism from BundleFuse's explicit
corrections, not evidence that considering several items jointly is new.

\section{State the contribution at the level the construction supports}
The proposed methodological contribution is the \emph{particular explicit scoring
construction for bundle completion}: historical relation estimates, category and
descriptive estimates, query-wise scale alignment, and specified optional corrections.
The paper must distinguish those design choices from the established ingredients.
A coherent combination can be useful, but naming its ingredients differently does
not establish novelty.

The central hypothesis is that carefully constructed relation scores provide a
competitive and inspectable alternative for this task. This is an affirmative,
testable hypothesis. It does not require asserting that neural models are excessive,
that all low-rank representations lose necessary information, or that the scorer
has a universal computational advantage.

A reviewer may reasonably ask whether BundleFuse is a feature-engineered ensemble.
The descriptive answer is that it combines engineered scores. The scientific answer
must specify the relationships those scores encode, why those relationships matter
for completion, and the distinctive behavior of the resulting estimator. A new
technical principle cannot be inferred solely from benchmark wins. The contribution
may ultimately be a well-supported construction and analysis rather than a general
new theory of recommendation.

The literature reviewed here does not establish that the complete combination is
unprecedented. A priority claim requires a closer component-by-component comparison
with the nearest published constructions and their underlying data. An explicit
score allows inspection of its arithmetic; it does not by itself establish causal
explanations or that the scores match human reasons.

\section{An introduction built from this path}
Completing a bundle requires selecting items that are appropriate companions to
those already present. Prior work on product relationships distinguishes
substitutability from complementarity, establishing that relevance to an item need
not mean resemblance~\cite{mcauley}. For completion, historical co-membership supplies
direct relational evidence, while item attributes allow that evidence to extend
beyond exact observed pairs.

Both explicit and neural models have established ways to exploit such information.
Shallow and side-information recommendation models demonstrate the viability of
explicit scoring~\cite{ease,sslim,fease}; modern bundle models learn representations
from several sources and relationships~\cite{clhe,ramen}. These developments motivate
examining the extent to which bundle compatibility can be estimated directly from
specified relations, with the contributions of those relations remaining visible.

We investigate this question through BundleFuse. The method constructs historical,
category and descriptive compatibility scores, places them on comparable scales,
and selects their combination using validation examples. It also provides explicit
candidate-frequency and multi-item corrections. This formulation makes the design
question concrete: which relational estimates are useful for completing a bundle,
and when do the corrections add information beyond their combination?

\paragraph{How this relates to the existing documents.}
This is the motivation and positioning that should precede
\texttt{publication\_argument\_technical.tex}. The historical manuscript's introduction
and related-work text contain stronger claims about cost, zero parameters, prior
methods and broad superiority that require reconciliation. They should not be used
as the current statement of the contribution. Experimental evidence belongs after
this argument and must remain within its verified scope.

\begin{thebibliography}{99}\small
\bibitem{hu} P. Hu et al. Hybrid Item-Item Recommendation via Semi-Parametric
Embedding. IJCAI, 2019. \url{https://doi.org/10.24963/ijcai.2019/350}.
\bibitem{coheat} H. Jeon et al. Cold-start Bundle Recommendation via
Popularity-based Coalescence and Curriculum Heating. WWW, 2024.
\url{https://doi.org/10.1145/3589334.3645377}.
\bibitem{gartrell} M. Gartrell, U. Paquet, N. Koenigstein. Low-Rank Factorization
of Determinantal Point Processes. AAAI, 2017.
\url{https://arxiv.org/abs/1602.05436}.
\bibitem{warlop} R. Warlop, J. Mary, M. Gartrell. Multi-Task Determinantal Point
Processes for Recommendation. Preprint, 2018.
\url{https://arxiv.org/abs/1805.09916}.
\bibitem{hoslim} E. Christakopoulou, G. Karypis. HOSLIM: Higher-Order Sparse
Linear Method for Top-N Recommender Systems. PAKDD, 2014.
\url{https://doi.org/10.1007/978-3-319-06605-9_4}.
\bibitem{breese} J. Breese, D. Heckerman, C. Kadie. Empirical Analysis of
Predictive Algorithms for Collaborative Filtering. UAI, 1998.
\url{https://arxiv.org/abs/1301.7363}.
\bibitem{kelen} D. Kelen et al. Efficient K-NN for Playlist Continuation.
RecSys Challenge, 2018. \url{https://doi.org/10.1145/3267471.3267477}.
\bibitem{bonnin} G. Bonnin, D. Jannach. A Comparison of Playlist Generation
Strategies for Music Recommendation and a New Baseline Scheme. 2013.
\url{https://web-ainf.aau.at/pub/jannach/files/Workshop_ITWP13.pdf}.
\bibitem{ludewig} M. Ludewig, D. Jannach. Evaluation of Session-based
Recommendation Algorithms. 2018. \url{https://arxiv.org/abs/1803.09587}.
\bibitem{zhou} D. Zhou, J. Huang, B. Sch\"olkopf. Learning with Hypergraphs:
Clustering, Classification, and Embedding. NeurIPS, 2006.
\url{https://papers.nips.cc/paper/2006/hash/dff8e9c2ac33381546d96deea9922999-Abstract.html}.
\bibitem{ferraro} A. Ferraro et al. Automatic Playlist Continuation Using a
Hybrid Recommender System Combining Features from Text and Audio. RecSys
Challenge, 2018. \url{https://doi.org/10.1145/3267471.3267473}.
\bibitem{sarwar} B. Sarwar et al. Item-Based Collaborative Filtering Recommendation
Algorithms. WWW, 2001. \url{https://www.ra.ethz.ch/CDstore/www10/papers/519/}.
\bibitem{melville} P. Melville, R. J. Mooney, R. Nagarajan. Content-Boosted
Collaborative Filtering for Improved Recommendations. AAAI, 2002.
\url{https://cdn.aaai.org/AAAI/2002/AAAI02-029.pdf}.
\bibitem{vasileva} M. Vasileva et al. Learning Type-Aware Embeddings for Fashion
Compatibility. ECCV, 2018. \url{https://arxiv.org/abs/1803.09196}.
\bibitem{volkovs} M. Volkovs et al. Two-stage Model for Automatic Playlist
Continuation at Scale. RecSys Challenge, 2018.
\url{https://doi.org/10.1145/3267471.3267480}.
\bibitem{deepicf} F. Xue et al. Deep Item-based Collaborative Filtering for Top-N
Recommendation. Preprint, 2018. \url{https://arxiv.org/abs/1811.04392}.

\bibitem{mcauley} J. McAuley, R. Pandey, J. Leskovec. Inferring Networks of Substitutable
and Complementary Products. KDD, 2015.
\url{https://cseweb.ucsd.edu/~jmcauley/pdfs/kdd15.pdf}.
\bibitem{ease} H. Steck. Embarrassingly Shallow Autoencoders for Sparse Data. WWW, 2019.
\url{https://arxiv.org/abs/1905.03375}.
\bibitem{sslim} X. Ning, G. Karypis. Sparse Linear Methods with Side Information for
Top-N Recommendations. WWW 2012 companion/poster version.
\url{https://archives.iw3c2.org/www2012/proceedings/companion/p581.pdf}.
\bibitem{fease} E. D. Cui, L. Zhang, W. P. Lee. Shallow AutoEncoding Recommender with
Cold Start Handling via Side Features. Preprint, 2025.
\url{https://arxiv.org/abs/2504.02288v4}.
\bibitem{clhe} Y. Ma et al. Leveraging Multimodal Features and Item-level User Feedback
for Bundle Construction. WSDM, 2024. \url{https://arxiv.org/abs/2310.18770}.
\bibitem{ramen} H.-S. Nguyen et al. RaMen: Multi-Strategy Multi-Modal Learning for Bundle
Construction. Preprint, 2025. \url{https://arxiv.org/abs/2507.14361}.
\bibitem{han} C. Han et al. Addressing Cold Start in Product Search via Empirical Bayes.
CIKM, 2022. \url{https://doi.org/10.1145/3511808.3557066}.
\end{thebibliography}
\end{document}
